

As model capabilities advance and they are deployed in increasingly
consequential settings, the need to monitor their internal state grows.
Their outputs cannot be taken at face value: reasoning models' stated
chains of thought are not always faithful to their internal reasoning
\citep{chen2025}. Detecting and removing undesirable behaviour, before and
during deployment, therefore needs tools that look inside the model.
\emph{Activation probes}, linear directions in the residual stream, are a
leading candidate: they flag a property at a fraction of the cost of an
auxiliary monitor model \citep{alain2018,mckenzie2025}, and a direction of
the same kind, added back into the stream, can steer the behaviour it
detects. Two questions decide whether a probe is useful: does it
\emph{detect} the property, in and out of distribution, and is its
direction \emph{causally implicated}, such that adding it changes what the
model does?

they usually test for either probing or steering, or 1 behaviour, or not every pooling technique, or not every estimator.
we wanted to study how all these design choices interact, and therefore built a unified pipeline
which studies different estimator techiques (), different pooling techniques (), different behaviours (), and for different purposes (probing/steering).
This is a comprehensive study, and makes observations mentioned in papers more robust:

like prior papers have found: 
Difference-of-means
(``mass-mean'') directions are reported as the most causally implicated for
truth despite comparable classification accuracy \citep{marks2024},
logistic-regression probes as the strongest detectors of strategic
deception \citep{goldowsky2025}, learned pooling as the best aggregation
for high-stakes detection \citep{mckenzie2025}, and mean-difference vectors
as effective steerers of open-ended generation \citep{rimsky2024}. 

This suggests things like a \emph{read/write dissociation}: the direction
that best reads a behaviour may not be the one that best writes it. Or that certain pooling schemes are better. 
But the
evidence comes from different models, behaviours, metrics and token
policies, so the dissociation could be an artefact of comparing them.

That's why we study all these design choices together. We also save the probes and use them to 
power an interactive interace for scoring ans steerng text directly, to democratize the inference process.

code, datasets, interface available at [todo]







===========

\section{Introduction}

As model capabilities advance and they are deployed in increasingly
consequential settings, the need to monitor their internal state grows.
Their outputs cannot be taken at face value: reasoning models' stated
chains of thought are not always faithful to their internal reasoning
\citep{chen2025}. Detecting and removing undesirable behaviour therefore requires tools that look inside the model.
\emph{Activation probes}, linear directions in the residual stream, are a
leading candidate: they flag a property at a fraction of the cost of an
auxiliary monitor model \citep{alain2018,mckenzie2025}, and a direction of
the same kind, added back into the stream, can steer the behaviour it
detects. Two questions decide whether a probe is useful: does it
\emph{detect} the property, in and out of distribution, and is its
direction \emph{causally implicated}, such that adding it changes what the
model does?

Existing studies typically examine one of these questions for a single behaviour, estimator or pooling scheme. We
wanted to see how these design choices interact, and so built one pipeline
that runs two estimators (difference-of-means and logistic regression),
three pooling schemes (last token, mean and learned attention) and three
behaviours (deception, high-stakes awareness and sycophancy) through the
same model, and evaluates every direction for both detection and steering.
Difference-of-means (``mass-mean'')
directions are reported as the most causally implicated for truth with lower 
classification accuracy \citep{marks2024}, logistic-regression
probes as the strongest detectors of strategic deception
\citep{goldowsky2025}, learned pooling as the best aggregation for
high-stakes detection \citep{mckenzie2025}, and difference-of-means vectors as
effective steerers of open-ended generation \citep{rimsky2024}.

Read together, these suggest a \emph{read/write dissociation}, in which
the direction that best reads a behaviour is not the one that best writes
it (logistic regression better for detection, difference-of-means better
for steering), and that some pooling schemes are simply better than
others. However, the evidence comes from different models, behaviours,
metrics and token policies, so either pattern could be an artefact of
comparing them. This is why we study these choices together: to make more robust observations about how estimator, pooling and behaviour determine a direction's effectiveness for detection and steering.


The exported probes and directions also power an interactive interface for
scoring and steering text directly (\S\ref{sec:impl}), to promote transparency and 
control in the inference process.

Code, datasets and
the interface are available at
\url{https://github.com/NimunB/probing-safety-behaviours}.
